\documentclass{beamer}
\usepackage{amsmath,amssymb,amsfonts}
\usepackage{algorithm}
\usepackage{algorithmic}
\usepackage{graphicx}
\usepackage{xcolor}

\usetheme{Madrid}
\usecolortheme{whale}

\title{Hybrid Active-Projection Method for Embedded QP Solving}
\author{NASOQ-MPC Project}
\date{\today}

\begin{document}

\begin{frame}
    \titlepage
\end{frame}

\begin{frame}{Outline}
    \tableofcontents
\end{frame}

\section{Problem Formulation}

\begin{frame}{Quadratic Programming (QP) Problem}
    Standard form QP:
    \begin{align}
        \min_x \quad & \frac{1}{2}x^T H x + q^T x \\
        \text{s.t.} \quad & Ax = b \\
        & Bx \leq d
    \end{align}
    
    \begin{itemize}
        \item $H \in \mathbb{R}^{n \times n}$ (symmetric positive definite)
        \item $q \in \mathbb{R}^n$
        \item $A \in \mathbb{R}^{m_e \times n}$ (equality constraints)
        \item $B \in \mathbb{R}^{m_i \times n}$ (inequality constraints)
    \end{itemize}
\end{frame}

\begin{frame}{Traditional Active-Set Methods}
    \begin{itemize}
        \item Maintain a set of active constraints: $\mathcal{A}(x) = \{i : B_i x = d_i\}$
        \item Iteratively update active set until optimality conditions are satisfied
        \item Each iteration:
            \begin{itemize}
                \item Solve equality-constrained QP with active constraints
                \item Add/remove one constraint at a time
                \item Compute step direction and length
                \item Update solution
            \end{itemize}
        \item Convergence is guaranteed but can be slow for large problems with many constraints
    \end{itemize}
\end{frame}

\section{Batch Constraint Processing}

\begin{frame}{Batch Processing of Constraints}
    \begin{itemize}
        \item \textbf{Traditional approach:} Add/remove one constraint at a time
        \item \textbf{Our approach:} Process multiple constraints in each iteration
    \end{itemize}
    
    \begin{theorem}[Batch Update]
        Given a feasible point $x_k$ and a set of violated constraints $\mathcal{V}_k$, we can compute a feasible step $\alpha_k d_k$ such that multiple constraints in $\mathcal{V}_k$ become active.
    \end{theorem}
\end{frame}

\begin{frame}{Batch Selection Algorithm}
    \begin{algorithm}[H]
        \caption{Batch Constraint Selection}
        \begin{algorithmic}[1]
            \STATE Compute violations: $v_i = B_i x_k - d_i$ for all $i \notin \mathcal{A}_k$
            \STATE Select batch $\mathcal{B}_k \subset \{i \notin \mathcal{A}_k : v_i > \text{tolerance}\}$
            \STATE Sort $\mathcal{B}_k$ by violation magnitude (decreasing)
            \STATE Keep top $b$ constraints in $\mathcal{B}_k$ (where $b$ is batch size)
        \end{algorithmic}
    \end{algorithm}
    
    \begin{itemize}
        \item \textbf{Theoretical Advantage:} Reduces the number of KKT systems to solve
        \item \textbf{Convergence:} Adding multiple constraints can lead to faster convergence
    \end{itemize}
\end{frame}

\begin{frame}{Optimal Step Length for Batch}
    For a batch $\mathcal{B}_k$ of constraints:
    
    \begin{align}
        \alpha_k = \min_{i \in \mathcal{B}_k} \frac{d_i - B_i x_k}{B_i d_k}
    \end{align}
    
    where $d_k$ is the search direction, and we only consider $i$ where $B_i d_k > 0$.
    
    \begin{theorem}[Batch Step Optimality]
        The optimal step length $\alpha_k$ ensures that at least one constraint in the batch becomes active without violating any constraint. With multiple constraining constraints, multiple constraints can become active in a single step.
    \end{theorem}
\end{frame}

\section{Matrix-Free KKT Solving}

\begin{frame}{KKT System for QP}
    At each iteration, we solve the KKT system:
    
    \begin{align}
        \begin{bmatrix}
            H & A^T & B_{\mathcal{A}}^T \\
            A & 0 & 0 \\
            B_{\mathcal{A}} & 0 & 0
        \end{bmatrix}
        \begin{bmatrix}
            d_k \\
            \lambda_e \\
            \lambda_i
        \end{bmatrix} =
        \begin{bmatrix}
            -\nabla f(x_k) \\
            0 \\
            0
        \end{bmatrix}
    \end{align}
    
    \begin{itemize}
        \item $B_{\mathcal{A}}$ contains rows of $B$ corresponding to active constraints
        \item $\lambda_e$ and $\lambda_i$ are Lagrange multipliers
    \end{itemize}
\end{frame}

\begin{frame}{Matrix-Free Approach}
    \begin{itemize}
        \item \textbf{Traditional approach:} Form and factorize KKT matrix
        \item \textbf{Our approach:} Only compute matrix-vector products
    \end{itemize}
    
    \begin{theorem}[Matrix-Free Iteration]
        Iterative methods like Conjugate Gradient (CG) and MINRES only require matrix-vector products of the form $Kv$ where $K$ is the KKT matrix and $v$ is a vector.
    \end{theorem}
\end{frame}

\begin{frame}{Matrix-Vector Products for KKT System}
    Given a vector $v = \begin{bmatrix} v_1 \\ v_2 \\ v_3 \end{bmatrix}$, the product $Kv$ is:
    
    \begin{align}
        Kv = 
        \begin{bmatrix}
            Hv_1 + A^Tv_2 + B_{\mathcal{A}}^Tv_3 \\
            Av_1 \\
            B_{\mathcal{A}}v_1
        \end{bmatrix}
    \end{align}
    
    Each component can be computed efficiently:
    \begin{itemize}
        \item $Hv_1$ using sparse matrix-vector product
        \item $A^Tv_2$ and $B_{\mathcal{A}}^Tv_3$ using sparse transpose operations
        \item $Av_1$ and $B_{\mathcal{A}}v_1$ using sparse matrix-vector products
    \end{itemize}
\end{frame}

\begin{frame}{Theoretical Advantages of Matrix-Free Approach}
    \begin{itemize}
        \item \textbf{Memory Efficiency:} 
            \begin{itemize}
                \item Memory complexity: $O(nnz(H) + nnz(A) + nnz(B))$ vs. $O(nnz(K))$
                \item KKT matrix can be much denser than original matrices
            \end{itemize}
        \item \textbf{Computational Efficiency:}
            \begin{itemize}
                \item Avoids expensive factorization: $O(n^3)$ for dense matrices
                \item Matrix-vector products: $O(nnz)$ operations
            \end{itemize}
        \item \textbf{Embedded Systems:}
            \begin{itemize}
                \item Predictable memory usage
                \item Avoids dynamic memory allocation
                \item Can exploit problem structure
            \end{itemize}
    \end{itemize}
\end{frame}

\section{Constraint Relaxation}

\begin{frame}{Constraint Relaxation Approach}
    \begin{itemize}
        \item \textbf{Traditional approach:} Enforce constraints strictly at each iteration
        \item \textbf{Our approach:} Relax constraints initially, then gradually tighten
    \end{itemize}
    
    Relaxed constraints:
    \begin{align}
        Bx \leq d + \gamma^{(k)} \cdot |d|
    \end{align}
    
    where $\gamma^{(k)}$ is the relaxation factor at iteration $k$.
\end{frame}

\begin{frame}{Relaxation Schedule}
    Adaptive relaxation schedule:
    \begin{align}
        \gamma^{(k+1)} = \max(\gamma^{(k)} \cdot \beta, \gamma_{\min})
    \end{align}
    
    \begin{itemize}
        \item $\gamma^{(0)}$: Initial relaxation factor (e.g., 0.1)
        \item $\beta$: Decay rate (e.g., 0.7)
        \item $\gamma_{\min}$: Minimum relaxation factor (e.g., 0.001)
    \end{itemize}
    
    \begin{theorem}[Constraint Tightening]
        If $\gamma^{(k)} \to 0$ as $k \to \infty$, then the solution converges to the solution of the original unrelaxed problem.
    \end{theorem}
\end{frame}

\begin{frame}{Theoretical Justification}
    \begin{itemize}
        \item \textbf{Improved Numerical Stability:}
            \begin{itemize}
                \item Avoids infeasible intermediate solutions
                \item Reduces zigzagging behavior
            \end{itemize}
        \item \textbf{Better Convergence Properties:}
            \begin{itemize}
                \item Similar to homotopy/continuation methods
                \item Can escape local minima in non-convex cases
            \end{itemize}
        \item \textbf{Faster Active-Set Identification:}
            \begin{itemize}
                \item Helps identify truly active constraints at optimality
                \item Reduces the need to add/remove constraints repeatedly
            \end{itemize}
    \end{itemize}
\end{frame}

\begin{frame}{End-to-End Approach}
    \begin{algorithm}[H]
        \caption{Hybrid Active-Projection Method}
        \begin{algorithmic}[1]
            \STATE Initialize $x_0$, $\mathcal{A}_0 = \emptyset$, $\gamma^{(0)}$
            \FOR{$k = 0, 1, \ldots$ until convergence}
                \STATE Apply relaxation: $\tilde{d} = d + \gamma^{(k)} \cdot |d|$
                \STATE Identify violated constraints batch $\mathcal{B}_k$
                \IF{$\mathcal{B}_k = \emptyset$}
                    \IF{$\gamma^{(k)} > \gamma_{\min}$}
                        \STATE Update $\gamma^{(k+1)} = \max(\gamma^{(k)} \cdot \beta, \gamma_{\min})$
                        \STATE Continue
                    \ELSE
                        \STATE Check feasibility with original constraints
                        \IF{feasible}
                            \STATE \textbf{return} $x_k$ (converged)
                        \ENDIF
                    \ENDIF
                \ENDIF
                \STATE Update active set: $\mathcal{A}_{k+1} = \mathcal{A}_k \cup \mathcal{B}_k$
                \STATE Solve KKT system matrix-free to get $d_k$
                \STATE Compute optimal step length $\alpha_k$
                \STATE Update: $x_{k+1} = x_k + \alpha_k d_k$
                \STATE Update relaxation: $\gamma^{(k+1)} = \max(\gamma^{(k)} \cdot \beta, \gamma_{\min})$
            \ENDFOR
        \end{algorithmic}
    \end{algorithm}
\end{frame}

\begin{frame}{Convergence Analysis}
    \begin{theorem}[Convergence Guarantee]
        If the original QP is feasible and bounded, the Hybrid Active-Projection Method converges to the optimal solution in a finite number of iterations.
    \end{theorem}
    
    \begin{proof}[Proof Sketch]
        \begin{itemize}
            \item Active set can change at most $2^{m_i}$ times
            \item Constraint relaxation ensures we approach feasibility
            \item Batch processing maintains descent property
            \item Matrix-free KKT solving preserves optimality conditions
        \end{itemize}
    \end{proof}
\end{frame}

\begin{frame}{Computational Complexity}
    \begin{itemize}
        \item \textbf{Per-iteration complexity:}
            \begin{itemize}
                \item Constraint violation check: $O(m_i \cdot nnz(B)/m_i) = O(nnz(B))$
                \item Batch selection: $O(m_i \log b)$ where $b$ is batch size
                \item Matrix-free KKT solve: $O(nnz(H) + nnz(A) + nnz(B)) \cdot O(\sqrt{\kappa})$
                  where $\kappa$ is condition number
            \end{itemize}
        \item \textbf{Number of iterations:} $O(m_i/b)$ compared to $O(m_i)$ for standard active-set
    \end{itemize}
\end{frame}

\section{Conclusion}

\begin{frame}{Summary}
    \begin{itemize}
        \item \textbf{Batch Constraint Processing:}
            \begin{itemize}
                \item Process multiple constraints simultaneously
                \item Reduces number of iterations
            \end{itemize}
        \item \textbf{Matrix-Free KKT Solving:}
            \begin{itemize}
                \item Avoids explicit KKT matrix formation
                \item Memory efficient for embedded systems
            \end{itemize}
        \item \textbf{Constraint Relaxation:}
            \begin{itemize}
                \item Start with relaxed problem, gradually tighten
                \item Improves numerical stability and convergence
            \end{itemize}
    \end{itemize}
    
    The combination of these techniques enables efficient QP solving on embedded systems with limited resources.
\end{frame}

\begin{frame}{Future Work}
    \begin{itemize}
        \item Adaptive batch sizing based on problem structure
        \item Specialized preconditioners for matrix-free KKT solving
        \item Warm-starting for sequential QP problems (e.g., MPC)
        \item Fixed-point arithmetic implementation for ultra-embedded systems
        \item Theoretical guarantees for non-convex problems
    \end{itemize}
\end{frame}

\begin{frame}
    \centering
    \Huge Thank you!
    
    \vspace{1cm}
    \large Questions?
\end{frame}

\end{document} 